% 原书物理页=印刷页3--8上；1.2全节.
\section{Lie 群}\label{sec:ssl-1-2}

\begin{definition}\label{def:ssl-1-2-1}
Lie 群是指一个光滑流形 $G$，其上配备一个光滑的群运算 $m\colon G\times G\longrightarrow G,\ (x,y)\mapsto xy$，并且取逆映射 $i\colon G\longrightarrow G,\ x\mapsto x^{-1}$ 也是光滑的.
\end{definition}

\begin{example*}
例如：任何离散群（姑且说是可数的）、加法群 $(\mathbb{R},+)$、乘法群 $(\mathbb{R}^{\times},\cdot)$、它的复类似 $(\mathbb{C}^{\times},\cdot)$，以及单位圆 $(S^{1},\cdot)$.
\end{example*}

\begin{numpara}[\newterm{ssl-general-linear-group}{一般线性群}{general linear group}]\label{np:ssl-1-2-2}
它们在某种意义上是 Lie 群的基本例子. 我们会在两种形式下遇到它们.

\textbf{矩阵形式.}\ $G=\mathrm{GL}_n(\mathbb{R})$，配备其作为 $M_n(\mathbb{R})\simeq\mathbb{R}^{n^{2}}$ 的开集的 $C^\infty$ 结构. 乘法 $(A,B)\mapsto AB$ 关于 $A,B$ 的表项是多项式的，故为 $C^\infty$. 逆映射\footnote{译注：原式中的 $\operatorname{com}(A)$ 未标转置，此处照录. 若 $\operatorname{com}(A)$ 表示代数余子式矩阵，则逆矩阵公式应再取转置.} $A\mapsto A^{-1}=\frac{1}{\det A}\,\operatorname{com}(A)$ 是解析的，从而特别地是 $C^\infty$ 的.

\textbf{内在形式.}\ 若 $V$ 是有限维 $\mathbb{R}$-向量空间，则令 $\mathrm{GL}(V):=\operatorname{Aut}_{\mathbb{R}}(V)$. 当然，若 $n$ 是 $V$ 的维数，则选取 $V$ 的一组基便给出一个双射 $\mathrm{GL}(V)\xrightarrow{\ \sim\ }\mathrm{GL}_n(\mathbb{R})$，它是下一定义意义下的 Lie 群同构. 同样，$\mathrm{GL}_n(\mathbb{C})$ 自然地是一个 Lie 群；且若 $V$ 是 $\mathbb{C}$-向量空间，则 $\mathrm{GL}(V):=\operatorname{Aut}_{\mathbb{C}}(V)$ 是一个 Lie 群.
\end{numpara}

\begin{definition}\label{def:ssl-1-2-3}
\newterm{ssl-lie-group-homomorphism}{Lie 群同态}{homomorphism of Lie groups}（有时我们干脆说光滑同态）是指一个映射 $\varphi\colon G\longrightarrow G'$，它是光滑的且与群运算交换.
\end{definition}

\begin{example*}
\begin{itemize}
\item 行列式 $\det\colon\mathrm{GL}_n(\mathbb{R})\longrightarrow\mathbb{R}^{\times}$.
\item 指数映射 $\exp\colon(\mathbb{R},+)\longrightarrow(\mathbb{R}_{>0},\cdot)$ 是一个同构，其逆为 $\log$.
\item 直线到圆周的缠绕映射：$\theta\in(\mathbb{R},+)\mapsto\exp(2i\pi\theta)\in S^{1}$ 是一个满射的\newterm{ssl-local-diffeomorphism}{局部微分同胚}{local diffeomorphism}.
\item 直线到环面的缠绕映射：$\theta\in(\mathbb{R},+)\mapsto(\exp(2i\pi\theta),\exp(2i\pi\theta\alpha))\in S^{1}\times S^{1}$. 若 $\alpha\in\mathbb{R}\setminus\mathbb{Q}$，则它是一个单的\newterm{ssl-immersion}{浸入}{immersion}，但不是\newterm{ssl-embedding}{嵌入}{embedding}，因为它的像是稠密的.
\end{itemize}
\end{example*}

\textbf{特别值得注意的特殊情形.}
\begin{itemize}
\item 形如 $\lambda\colon(\mathbb{R},+)\longrightarrow G$ 的同态（不甚严格地）称为\newterm{ssl-one-parameter-subgroup}{单参数子群}{one-parameter subgroup}.
\item 形如 $\rho\colon G\longrightarrow\mathrm{GL}(V)$ 的同态称为 $G$ 在 $V$ 上的（线性）表示. 空间 $V$ 也称为表示 $\rho$ 的底层空间.
\end{itemize}

\begin{numpara}[全纯变体]\label{np:ssl-1-2-4}
\newterm{ssl-complex-lie-group}{复 Lie 群}{complex Lie group}是指一个复流形，其上带有群结构，使得乘法与取逆都是全纯的. 复 Lie 群之间的态射就是全纯的群同态. 例如，$\mathrm{GL}_n(\mathbb{C})$ 是复 Lie 群，且指数映射 $(\mathbb{C},+)\longrightarrow(\mathbb{C}^{\times},\cdot)$ 是全纯同态.
\end{numpara}

\begin{numpara}[代数变体]\label{np:ssl-1-2-5}
设 $\mathbb{K}$ 是\newterm{ssl-algebraically-closed-field}{代数闭域}{algebraically closed field}. $\mathbb{K}$ 上的\newterm{ssl-affine-algebraic-group}{仿射代数群}{affine algebraic group}是指 $\mathbb{K}^{n}$ 的一个由多项式方程定义的子集，其上带有群结构，且乘积与取逆都是坐标的多项式. 例如 $\mathrm{GL}_n(\mathbb{K})$ 是 $\mathbb{K}$ 上的代数群. 反之，指数映射一般没有代数意义上的类似物.
\end{numpara}

\begin{definition}\label{def:ssl-1-2-6}
Lie 群 $G$ 的\newterm{ssl-lie-subgroup}{Lie 子群}{Lie subgroup} $H$ 是指一个同时又是子群的\newterm{ssl-submanifold}{子流形}{submanifold}. 等价地说，单射 $H\hookrightarrow G$ 是一个 Lie 群的嵌入.
\end{definition}

\begin{example*}
\begin{itemize}
\item $G=\mathbb{R}$ 中的 $H=\mathbb{Z}$.
\item $G=\mathbb{C}^{\times}$ 中的 $H=S^{1}$.
\end{itemize}
\end{example*}

由于是子流形，Lie 子群是局部闭的. 利用平移作用容易证明它实际上是闭的. 下面的定理给出了一个颇为惊人的逆命题，它是大量例子的来源.

\begin{theorem}[Cartan，Von-Neumann]\label{thm:ssl-1-2-7}
Lie 群 $G$ 的任何闭子群 $H$ 都是 Lie 子群. 更精确地说，$H$ 上存在（必然唯一的）Lie 群结构，使 $H$ 成为 $G$ 的子流形.
\end{theorem}

我们把证明交给讲义 [GALM2]\footnote{译注：[GALM2] 指本系列另一册《Lie 群与 Lie 代数》，编号为其节内统一编号.}. 这个定理的实际价值在于：由此很容易产生大量 Lie 群的新例子. 下述情形尤其多产：设 $G$ 光滑地作用在光滑流形 $M$ 上（即作用 $G\times M\longrightarrow M$ 是光滑的）. 那么，一点 $m\in M$ 的\newterm{ssl-stabilizer}{稳定子}{stabilizer} $G_m$ 是闭的，从而是一个 Lie 群.

特别地，所有\newterm{ssl-classical-group}{经典群}{classical group}都是 Lie 群. 由于它们是最有意思的例子，我们在下面回顾其清单. 字母 $\mathbb{K}$ 不加区分地表示域 $\mathbb{R}$ 与 $\mathbb{C}$，而 $V$ 表示一个维数为 $n$ 的有限维 $\mathbb{K}$-向量空间.

\begin{numpara}[特殊线性群]\label{np:ssl-1-2-8}
行列式 $\mathrm{GL}(V)\longrightarrow\mathrm{GL}_1(\mathbb{K})$ 是拓扑群的连续同态，故其核 $\mathrm{SL}(V)$ 是闭的，从而是 Lie 子群. 我们称它为 $V$ 的\newterm{ssl-special-linear-group}{特殊线性群}{special linear group}. 选取 $V$ 的一组基便可把它等同于 $\mathrm{SL}_n(\mathbb{K})$.
\end{numpara}

\begin{exercise*}
$\mathrm{SL}_n(\mathbb{K})$ 的中心由行列式为 $1$ 的位似变换组成，因此它同构于 $\mathbb{K}^{\times}$ 中 $n$ 次单位根所成的群 $\mu_n(\mathbb{K})$（当 $\mathbb{K}=\mathbb{C}$ 时为 $n$ 阶，当 $\mathbb{K}=\mathbb{R}$ 时为 $1$ 阶或 $2$ 阶）.
\end{exercise*}

\begin{numpara}[正交群]\label{np:ssl-1-2-9}
记 $\operatorname{Sym}_{\mathbb{K}}(V)$ 为 $\mathbb{K}$ 上的\newterm{ssl-symmetric-bilinear-form}{对称双线性型}{symmetric bilinear form} $\Phi\colon V\times V\longrightarrow\mathbb{K}$ 所成的 $\mathbb{K}$-向量空间. $\mathrm{GL}(V)$ 在 $\operatorname{Sym}_{\mathbb{K}}(V)$ 上有一个线性从而连续的作用，由 $g\Phi(v,w):=\Phi(g^{-1}v,g^{-1}w)$ 给出. 于是 $\Phi$ 在此作用下的稳定子是 $\mathrm{GL}(V)$ 中的闭集，从而是 Lie 子群. 把它记作
\[
\mathrm{O}(\Phi):=\{g\in\mathrm{GL}(V),\ \Phi(gv,gw)=\Phi(v,w),\ \forall v,w\in V\}.
\]
若 $\Phi$ 非退化，则称之为 $\Phi$ 的\newterm{ssl-orthogonal-group}{正交群}{orthogonal group}.

我们回顾：两个\newterm{ssl-quadratic-form}{二次型}{quadratic form}称为等价的，如果它们属于 $\mathrm{GL}(V)$ 作用下的同一轨道. 由此可知，它们的正交群在 $\mathrm{GL}(V)$ 中互相共轭，特别地互相同构. 当 $\mathbb{K}=\mathbb{C}$ 时，我们知道所有二次型都是等价的，于是 $\mathrm{O}(\Phi)$ 同构于通常的复正交群
\[
\mathrm{O}(n,\mathbb{C})=\{M\in M_n(\mathbb{C}),\ {}^{t}M\cdot M=I_n\}.
\]
当 $\mathbb{K}=\mathbb{R}$ 时，Sylvester 定理断言存在一个整数 $p\leqslant n$ 以及一组基 $e_1,\cdots,e_n$，使得在其中
\[
\Phi\Big(\sum_{i=1}^{n}x_i e_i,\ \sum_{i=1}^{n}y_i e_i\Big)=\sum_{i=1}^{p}x_i y_i-\sum_{i=p+1}^{n}x_i y_i.
\]
整数对 $(p,\ q:=n-p)$ 称为 $\Phi$ 的\newterm{ssl-signature}{惯性指数对}{signature}. 于是 $\mathrm{O}(\Phi)$ 同构于矩阵群
\[
\mathrm{O}(p,q):=\{M\in M_{p+q}(\mathbb{R}),\ {}^{t}M\cdot D_{p,q}\cdot M=D_{p,q}\},
\]
其中 $D_{p,q}=D(1,\cdots,1,-1,\cdots,-1)$ 是对角矩阵，$1$ 重复 $p$ 次，$-1$ 重复 $q$ 次. 当 $p=n$ 与 $p=0$ 时，我们重新得到 Euclid 正交群 $\mathrm{O}(n,0)=\mathrm{O}(n)$，众所周知它是紧的. 其余的正交群都不是紧的. 我们有 $\mathrm{O}(p,q)=\mathrm{O}(q,p)$，并且可以证明诸 $\mathrm{O}(p,q)$（$q\leqslant p$）两两不同构.
\end{numpara}

\begin{example*}
对于时空 $\mathbb{R}^{4}$ 上的型 $xx'+yy'+zz'-tt'$，如此得到的是 Lorentz 群 $\mathrm{O}(3,1)$，它出现于广义相对论中.
\end{example*}

另外，我们用下式定义特殊正交群：
\[
\mathrm{SO}(p,q):=\mathrm{O}(p,q)\cap\mathrm{SL}_{p+q}(\mathbb{R}).
\]
\begin{numpara}[酉群]\label{np:ssl-1-2-10}
这里 $\mathbb{K}=\mathbb{C}$. 记 $\operatorname{Herm}(V)$ 为\newterm{ssl-hermitian-form}{Hermite 型}{hermitian form} $\Psi\colon V\times V\longrightarrow\mathbb{C}$ 所成的 $\mathbb{R}$-向量空间. $\mathrm{GL}(V)$ 在 $\operatorname{Herm}(V)$ 上有一个线性从而连续的作用，由 $g\Psi(v,w):=\Psi(g^{-1}v,g^{-1}w)$ 给出. 于是 $\Psi$ 在此作用下的稳定子是 $\mathrm{GL}(V)$ 中的闭集，从而是 Lie 子群. 把它记作
\[
\mathrm{U}(\Psi):=\{g\in\mathrm{GL}(V/\mathbb{C}),\ \Psi(gv,gw)=\Psi(v,w),\ \forall v,w\in V\}.
\]
若 $\Psi$ 非退化，则称之为 $\Psi$ 的\newterm{ssl-unitary-group}{酉群}{unitary group}.

与对称情形一样，两个 Hermite 型称为等价的，如果它们属于 $\mathrm{GL}(V)$ 作用下的同一轨道. 这时它们的酉群互相同构，甚至互相共轭. 这里同样地，Hermite 型按其惯性指数对分类：Sylvester 定理断言存在一个整数 $p\leqslant n$ 以及一组基 $e_1,\cdots,e_n$，使得在其中
\[
\Psi\Big(\sum_{i=1}^{n}x_i e_i,\ \sum_{i=1}^{n}y_i e_i\Big)=\sum_{i=1}^{p}x_i\bar y_i-\sum_{i=p+1}^{n}x_i\bar y_i.
\]
于是，置 $q=n-p$，拓扑群 $\mathrm{U}(\Psi)$ 同构于矩阵群
\[
\mathrm{U}(p,q):=\{M\in M_{p+q}(\mathbb{C}),\ {}^{t}\bar M\cdot D_{p,q}\cdot M=D_{p,q}\}.
\]
当 $pq=0$ 时，我们重新得到通常的酉群 $\mathrm{U}(n,0)=\mathrm{U}(n)$，众所周知它是紧的. 其余的酉群都不是紧的. 我们有 $\mathrm{U}(p,q)=\mathrm{U}(q,p)$，并且可以证明诸 $\mathrm{U}(p,q)$（$q\leqslant p$）两两不同构.
\end{numpara}

\begin{remark*}
群 $\mathrm{U}(1,n-1)$ 在当今数论的诸多问题中起着举足轻重的作用.
\end{remark*}

与前面一样，我们有特殊酉群
\[
\mathrm{SU}(p,q):=\mathrm{U}(p,q)\cap\mathrm{SL}_{p+q}(\mathbb{C}).
\]
\begin{exercise*}
求出 $\mathrm{U}(n)$、$\mathrm{SU}(n)$、$\mathrm{O}(n)$、$\mathrm{SO}(n)$ 的中心.
\end{exercise*}

\begin{numpara}[辛群]\label{np:ssl-1-2-11}
再次设 $\mathbb{K}$ 表示 $\mathbb{R}$ 或 $\mathbb{C}$. 记 $\operatorname{Alt}_{\mathbb{K}}(V)$ 为 $\mathbb{K}$ 上的\newterm{ssl-alternating-bilinear-form}{交错双线性型}{alternating bilinear form} $\psi\colon V\times V\longrightarrow\mathbb{K}$ 所成的 $\mathbb{K}$-向量空间. $\mathrm{GL}(V)$ 在 $\operatorname{Alt}_{\mathbb{K}}(V)$ 上有一个线性从而连续的作用，由 $g\psi(v,w):=\psi(g^{-1}v,g^{-1}w)$ 给出. 于是 $\psi$ 在此作用下的稳定子是 $\mathrm{GL}(V)$ 中的闭集，从而是 Lie 子群. 把它记作
\[
\mathrm{Sp}(\psi):=\{g\in\mathrm{GL}(V),\ \psi(gv,gw)=\psi(v,w),\ \forall v,w\in V\}.
\]
若 $\psi$ 非退化（这时称它为\newterm{ssl-symplectic-form}{辛形式}{symplectic form}），则称之为 $\psi$ 的\newterm{ssl-symplectic-group}{辛群}{symplectic group}. 此时 $V$ 的维数必然是偶数，我们把记号改设为 $\dim_{\mathbb{K}}V=2n$. 与前面一样，这个群在同构意义下只依赖于 $\psi$ 的等价类. 而我们知道，所有辛形式都是等价的. 特别地，存在 $V$ 的一组基 $e_1,\cdots,e_{2n}$，在其中\footnote{译注：下一式左边两处求和上限原文均为 $n$，此处照录；按 $2n$ 维空间的坐标展开，应为 $2n$.}
\[
\psi\Big(\sum_{i=1}^{n}x_i e_i,\ \sum_{i=1}^{n}y_i e_i\Big)=\sum_{i=1}^{n}(x_i y_{n+i}-x_{n+i}y_i).
\]
换句话说，$\psi$ 在这组基下的矩阵是反对称矩阵 $J=J_{2n}=\begin{pmatrix}0_n&I_n\\-I_n&0_n\end{pmatrix}$. 于是可见 $\mathrm{Sp}(\psi)$ 同构于辛群
\[
\mathrm{Sp}_{2n}(\mathbb{K})=\{M\in\mathrm{GL}_{2n}(\mathbb{K}),\ {}^{t}MJM=J\}.
\]
由此，$\mathrm{Sp}_{2}(\mathbb{K})=\mathrm{SL}_{2}(\mathbb{K})$. 更一般地，可以证明辛矩阵的行列式总等于 $1$，因此对一切 $n$，$\mathrm{Sp}_{2n}(\mathbb{K})$ 都是 $\mathrm{SL}_{2n}(\mathbb{K})$ 的闭的非紧子群.
\end{numpara}

\begin{exercise*}
证明：$\mathbb{C}$-向量空间 $V$ 上一个 Hermite 型的实部（相应地，虚部）是底层 $\mathbb{R}$-向量空间 $\widetilde{V}$ 上的对称（相应地，交错）$\mathbb{R}$-双线性型. 由此推出：在 $\mathrm{GL}_n(\mathbb{C})\simeq\{P\in\mathrm{GL}_{2n}(\mathbb{R}),\ P^{-1}J_{2n}P=J_{2n}\}\subset\mathrm{GL}_{2n}(\mathbb{R})$ 中，有 $\mathrm{U}(n)=\mathrm{O}(2n)\cap\mathrm{Sp}_{2n}(\mathbb{R})$.
\end{exercise*}

\begin{numpara}[经典群的重要性]\label{np:ssl-1-2-12}
它们出现在相当多样的领域中.

\textbf{几何.}\ 这是它们最得心应手的领域. 这些群常常是一些著名空间（球面、双曲面……）的对称群，并且正因如此，它们是非欧几何的模型.

\textbf{数论.}\ 这是一个稍显出乎意料的应用，但近百年来，人们在形如 $G(\mathbb{Z})\backslash G(\mathbb{R})/K$ 的空间上研究函数空间、微分形式等，其中 $G$ 是实经典群，$K$ 是 $G$ 的一个极大紧子群. 例如，$G=\mathrm{GL}_1$ 时得到三角函数，$G=\mathrm{GL}_2$ 时得到模形式，对一般的 $G$ 则说自守形式. 这些对象起着至关重要的作用，例如在 Fermat 定理的证明中.

\textbf{物理.}
\begin{itemize}
\item 广义相对论的 Lorentz 群是 $\mathrm{O}(1,3)$. 它是与时空几何相联系的群.
\item 某些基本粒子由 $\mathrm{SU}(3)$ 的\newterm{ssl-irreducible-representation}{不可约表示}{irreducible representation}分类.
\item 氢原子的能级由球面 $S^{2}$ 上球面 Laplace 算子的特征值给出，我们稍后将利用 $C^{\infty}(S^{2})$ 在 $\mathrm{SO}_3(\mathbb{R})$ 作用下的谱分解来计算它们.
\end{itemize}

\begin{remark*}
群 $\mathrm{SL}_n(\mathbb{C})$、$\mathrm{Sp}_{2n}(\mathbb{C})$ 与 $\mathrm{O}_n(\mathbb{C})$ 都是复 Lie 群，甚至是代数群. 群 $\mathrm{U}(n)$ 不是复 Lie 群（但它由 $\mathbb{R}$ 上的多项式映射定义）.
\end{remark*}
\end{numpara}
